\chapter{Uso responsable de la IA}
\label{anexo:uso-ia}
Este anexo declara el apoyo de inteligencia artificial documentado durante la preparación de la memoria. Su alcance comprende las tareas observadas en el trabajo con el asistente; no constituye un inventario de otras herramientas que el autor haya podido emplear fuera de estas sesiones.

\section{Herramientas de IA utilizadas}
Se ha utilizado el asistente de OpenAI integrado en Codex como apoyo para consultar el proyecto local, revisar código e historial, proponer redacción y editar los archivos LaTeX de Overleaf. También se ha empleado como ayuda para preparar y corregir fuentes PlantUML y explicaciones de diagramas. No se especifica una versión de modelo que no conste en los registros disponibles.

\section{Para qué se ha usado y para qué no}
La asistencia se ha utilizado para organizar información técnica, relacionar requisitos con PBI, redactar explicaciones de arquitectura y completar los anexos respetando su estructura. En los diagramas se ha empleado para representar despliegue, acceso, solicitud de recogida y estados del transporte y para ajustar las explicaciones a la revisión del código consultada.

En esta actualización se han contrastado las capturas y los datos ya recogidos en la memoria con el repositorio. Se ha conservado la tabla de horas del Sprint 1 y se ha diferenciado la implementación actual de las evidencias históricas. Los requisitos del Sprint 2 se actualizan al cierre de su implementación y se conserva la separación entre evidencias históricas y resultados actuales. La asistencia también se utilizó para implementar y probar P23 y ampliar P24, incluidos los reintentos, transacciones y pruebas con datos ficticios; las reglas de estado fueron confirmadas por el autor.

La IA no se considera fuente primaria de fechas, dedicación, estados de DoneTonic ni resultados de pruebas. No se han usado estimaciones generadas por el asistente para completar horas reales o fechas ausentes. Tampoco se presenta una descripción generada de una función como prueba de que esa función existe, ni una prueba simulada como evidencia de funcionamiento real de Google o del correo.

\section{Cómo se ha verificado lo generado}
La revisión de contenidos se apoya en tres tipos de evidencia: archivos de la aplicación para describir sus capacidades actuales; historial Git para situar cambios técnicos; y registros de planificación aportados para las estimaciones y el seguimiento. Estas fuentes responden a preguntas distintas: la fecha de un commit no fija por sí sola el inicio o cierre de un sprint.

Un ejemplo de corrección fue la persistencia de sesiones. Una descripción anterior indicaba que se perdían al reiniciar el servidor. Al contrastar una revisión posterior de \texttt{sesiones.js} y la capa de datos, se actualizó el diagrama y su explicación para reflejar sesiones persistentes con caducidad de ocho horas. La corrección ilustra la necesidad de comprobar las propuestas del asistente frente al código vigente.

Las pruebas automatizadas existentes constituyen un mecanismo adicional de validación técnica cuando se ejecutan y se revisan sus resultados. La consulta de sus archivos no equivale a una ejecución satisfactoria. La ejecución CI 44078563639 verifica 85 casos y se completa con una prueba visual local del transporte. La revisión documental no certifica pruebas reales de proveedores externos. Las discrepancias de los gráficos de DoneTonic se conservan y explican, sin completar ni inventar los datos ausentes. La compilación de LaTeX y la revisión del PDF comprueban la integridad y presentación del documento, sin sustituir la validación funcional de la aplicación.

\section{Autoría y responsabilidad}
El autor mantiene la responsabilidad sobre las decisiones de diseño, la exactitud de las afirmaciones y la versión entregada. La asistencia automatizada aporta propuestas y realiza ediciones, pero no sustituye la comprensión del proyecto ni la revisión crítica de los resultados.

Antes de la entrega, reviso todos los requisitos, confirmo los datos de seguimiento y completo las comprobaciones pendientes. Las afirmaciones deben poder justificarse mediante código, registros o pruebas identificables. Cuando una evidencia no está disponible, se explicita su alcance en lugar de atribuir al proyecto resultados no demostrados.
